\Page{Défis avancés\quad /\quad D1}{Combien de plans sont nets ?}{defis1}
\phantomsection\label{ex-D1}
\source{Problème original OptiTAB. Niveau 3 ; 45 minutes. Indices p.\,\pageref{indices} ; solution p.\,\pageref{corr-d1}.}
Un objectif est modélisé par une lentille de focale $f'>0$ et un diaphragme de diamètre $a=f'/N$ dans son plan. Il est réglé sur un objet à distance $p_0>f'$ ; le capteur est fixé à $q_0=f'p_0/(p_0-f')$. Un point est acceptable si sa tache de flou a un diamètre inférieur à $\varepsilon$.
\fig{profondeur}{Reprendre la construction en précisant les deux plans de convergence et le plan du capteur.}
\begin{questions}
\item Retrouver $b=a|1-q_0/q|$ par similitude. Montrer ensuite que
\[\left|\frac1p-\frac1{p_0}\right|\leq k,\qquad k=\frac{\varepsilon}{a q_0}.\]
\item En déduire la borne proche $p_-$ et la borne lointaine $p_+$ de la zone acceptable. Distinguer les cas $k<1/p_0$ et $k\geq1/p_0$.
\item Calculer les deux bornes pour $f'=50{,}0\ \mathrm{mm}$, $p_0=5{,}00\ \mathrm{m}$, $N=8{,}00$ et $\varepsilon=30{,}0\ \mu\mathrm{m}$.
\item Déterminer la distance de réglage $H$ à partir de laquelle la borne lointaine atteint l'infini. Vérifier que $H=f'+f'^2/(N\varepsilon)$ dans ce modèle.
\item Expliquer ce que ce modèle ne permet pas de conclure lorsque l'on ferme fortement le diaphragme.
\end{questions}
\begin{methode}{Ce qui rend ce défi formateur}
Une valeur absolue cache deux inégalités. Un dénominateur nul ou négatif change le sens physique d'une borne ; il faut le discuter au lieu d'appliquer une expression aveuglément.
\end{methode}

\Page{Défis avancés\quad /\quad D2}{Dimensionner une liaison optique}{defis2}
\phantomsection\label{ex-D2}
\source{Problème original OptiTAB. Niveau 3 ; 40 minutes. Aucune connaissance de télécommunications n'est supposée.}
Une liaison utilise une fibre à saut d'indice $n_1=1{,}48$, $n_2=1{,}46$. Une impulsion injectée a une durée initiale $\tau_0=20{,}0\ \mathrm{ns}$. Pour ce modèle d'ingénierie simplifié, on définit la durée reçue par $\tau_{\rm rec}=\tau_0+\Delta t$ et l'on impose $2\tau_{\rm rec}\leq1/R$, où $R$ est le débit binaire.
\fig{dispersion}{Le temps de transit et la dispersion intermodale ne jouent pas le même rôle dans le cahier des charges.}
\begin{questions}
\item Établir l'expression exacte de $\Delta t$ en fonction de $L$, puis écrire $\Delta t=\kappa L$ en donnant l'unité de $\kappa$.
\item Une fibre de $1{,}00\ \mathrm{km}$ convient-elle à $R=10{,}0\ \mathrm{Mbit\,s^{-1}}$ selon le critère fourni ?
\item Déterminer la longueur maximale acceptable. Quelle impossibilité apparaît si $1/(2R)\leq\tau_0$ ?
\item Dans l'air, un connecteur injecte un rayon méridien à $12{,}0^\circ$ de l'axe. Vérifier son guidage ; déterminer l'angle interne et le temps de transit à travers la longueur maximale calculée.
\item On propose de réduire $n_1-n_2$ en conservant $n_1$ fixé. Prévoir qualitativement l'effet sur l'étalement et sur le cône d'entrée. Quel compromis apparaît ?
\end{questions}
\begin{docu}{Ce qui est donné, ce qui est à établir}
La règle d'addition des durées et la marge de séparation 2 sont imposées pour ce problème. La dépendance géométrique de $\Delta t$ et la condition de guidage sont à démontrer. On ne tient compte ni des pertes ni de la dispersion chromatique.
\end{docu}

\Page{Défis avancés\quad /\quad D3}{Un microscope à deux lentilles}{defis3}
\phantomsection\label{ex-D3}
\source{Problème original OptiTAB. Niveau 3 ; 50 minutes. Prolongement : sujet Mines-Ponts MP 2021, parcours p.\,\pageref{auth-concours}.}
Un microscope comporte un objectif de focale $f'_1=4{,}00\ \mathrm{mm}$ et un oculaire de focale $f'_2=25{,}0\ \mathrm{mm}$. L'intervalle optique est $\Delta=\overline{F'_1F_2}=160\ \mathrm{mm}$. L'image finale est à l'infini ; les deux lentilles sont dans l'air et dans les conditions de Gauss.
\fig{microscope}{Schéma de principe calculé avec des proportions lisibles, différentes des valeurs du problème ; les rayons émergents sont parallèles.}
\begin{questions}
\item Où doit se trouver l'image intermédiaire ? Déterminer $\overline{O_1A_1}$ et la distance entre les lentilles.
\item Calculer la position de l'objet par rapport à $O_1$, puis sa distance au foyer $F_1$.
\item Déterminer le grandissement de l'objectif. Pour un objet de hauteur $h=10{,}0\ \mu\mathrm{m}$, calculer la hauteur de l'image intermédiaire.
\item Définir $\alpha'$ comme la valeur absolue du petit angle sous lequel l'image finale est observée. Établir la puissance intrinsèque $P_i=\alpha'/h=\Delta/(f'_1f'_2)$ et la calculer.
\item Comparer l'angle obtenu à celui de la vision directe du même objet à $d_m=25{,}0\ \mathrm{cm}$. Donner le grossissement angulaire en valeur absolue.
\item Expliquer l'intérêt de l'image finale à l'infini pour un œil emmétrope. Est-ce une preuve que tous les détails de l'objet sont résolus ?
\end{questions}

\Page{Évaluation\quad /\quad DM}{Mesurer une focale par Bessel}{dm}
\phantomsection\label{ex-DM}
\source{Devoir maison original OptiTAB. Travail estimé : 1 h 30 ; calculatrice autorisée. Barème pédagogique sur 20. À rendre avec un schéma et une discussion des incertitudes.}
On place un objet et un écran à une distance fixe $D$. Une lentille convergente peut coulisser entre eux. On observe deux positions de netteté séparées de $d$. Les distances sont mesurées indépendamment :
\[D=(1000{,}0\ \mathrm{mm}),\quad u(D)=1{,}0\ \mathrm{mm},\qquad
d=(600{,}0\ \mathrm{mm}),\quad u(d)=2{,}0\ \mathrm{mm}.\]
\fig{bessel}{Les positions sont mesurées par rapport à l'objet. Les schémas et la notation doivent rester cohérents.}
\begin{questions}
\item En notant $x$ la distance objet-lentille, établir l'équation donnant les positions nettes et la condition sur $D$ et $f'$. \textbf{4 points}
\item Exprimer les deux solutions et montrer que $f'=(D^2-d^2)/(4D)$. \textbf{3 points}
\item Calculer $f'$ et les deux positions. \textbf{2 points}
\item Déterminer les deux grandissements et interpréter leur produit. \textbf{2 points}
\item Calculer $u(f')$ par propagation au premier ordre avec $D,d$ indépendants. Présenter le résultat. \textbf{3 points}
\item Une méthode indépendante donne $f'_r=158{,}5\ \mathrm{mm}$ et $u(f'_r)=0{,}8\ \mathrm{mm}$. Calculer l'écart normalisé et commenter. \textbf{2 points}
\item Proposer une méthode pour estimer l'incertitude de mise au point. \textbf{2 points}
\item Citer deux limites physiques ou expérimentales susceptibles de biaiser le résultat. \textbf{2 points}
\end{questions}

\Page{Évaluation\quad /\quad DS\quad /\quad 1 sur 2}{Interfaces et projection}{ds1}
\phantomsection\label{ex-DS}
\source{Devoir surveillé original OptiTAB. Durée : 2 h ; calculatrice autorisée ; note sur 20. Les deux pages constituent un seul sujet. Toute valeur doit être justifiée par une expression littérale.}
\Sub{Partie A\quad /\quad Un capteur par réflexion totale\quad /\quad 8 points}
Une source envoie un rayon dans une plaque de verre d'indice $n_v=1{,}50$. L'angle dans le verre avec la normale à la face extérieure vaut $i=50{,}0^\circ$. Cette face est en contact avec l'air ($n_a=1{,}00$) ou avec une goutte d'eau ($n_e=1{,}33$). Un détecteur reçoit le faisceau réfléchi. Le trajet d'injection n'est pas étudié.
\fig{capteur-pluie}{Le dispositif compare la lumière réfléchie quand la face extérieure est sèche ou mouillée.}
\begin{questions}
\item Énoncer les lois de réflexion et de réfraction en identifiant les angles et les indices. \textbf{2 points}
\item Établir et calculer les deux angles critiques. \textbf{2 points}
\item Décider du comportement à $50{,}0^\circ$ dans chaque situation et calculer l'angle transmis quand il existe. \textbf{2 points}
\item Donner l'intervalle d'incidences permettant de distinguer les deux situations par réflexion totale ; expliquer la diminution du signal réfléchi en présence d'eau. \textbf{2 points}
\end{questions}
\Sub{Partie B\quad /\quad La projection d'un repère\quad /\quad 8 points}
Une lentille de focale $f'=+100\ \mathrm{mm}$ reçoit la lumière d'un objet réel de hauteur $h=12{,}0\ \mathrm{mm}$ placé à $300\ \mathrm{mm}$ à sa gauche. L'axe est orienté de gauche à droite. Les questions de cette partie se poursuivent page suivante.
\begin{questions}
\item Écrire la distance objet algébrique et la relation de conjugaison adaptée. \textbf{1 point}
\end{questions}

\Page{Évaluation\quad /\quad DS\quad /\quad 2 sur 2}{Mesurer et transmettre}{ds2}
\Sub{Partie B\quad /\quad Suite}
\begin{enumerate}[label=\textcolor{otNavy}{\bfseries B\arabic*.},start=2]
\item Déterminer la position de l'image, son grandissement et sa hauteur algébrique. Donner sa nature et son orientation. \textbf{3 points}
\item On fixe ensuite la distance objet-écran à $D=450\ \mathrm{mm}$. Vérifier la condition de projection et déterminer les deux positions de netteté de la même lentille. \textbf{2 points}
\item Réaliser une construction justifiée à deux rayons pour la position initiale. Que deviennent le grandissement et la taille de l'image pour l'autre position à $D=450\ \mathrm{mm}$ ? \textbf{2 points}
\end{enumerate}
\Sub{Partie C\quad /\quad Une fibre de 500 mètres\quad /\quad 4 points}
Une fibre à saut d'indice a un cœur d'indice $n_1=1{,}48$, une gaine d'indice $n_2=1{,}46$ et une longueur $L=500\ \mathrm{m}$. On étudie les rayons méridiens et on néglige la dispersion matérielle. Le milieu d'entrée est l'air. On prend $c=3{,}00\times10^8\ \mathrm{m\,s^{-1}}$.
\begin{enumerate}[label=\textcolor{otNavy}{\bfseries C\arabic*.}]
\item À partir de la condition à la paroi, établir l'angle d'acceptance et le calculer. \textbf{2 points}
\item Établir l'étalement maximal $\Delta t$, puis estimer le débit maximal avec le critère explicite $T\geq2\Delta t$. Les impulsions injectées sont supposées de durée négligeable. \textbf{2 points}
\end{enumerate}
\begin{docu}{Gestion du temps proposée}
Lecture et schémas : 10 min ; partie A : 30 min ; partie B : 35 min ; partie C : 30 min ; relecture : 15 min. Un résultat incohérent doit être signalé et discuté plutôt que masqué.
\end{docu}
\begin{methode}{Autoévaluation avant le corrigé}
\checkline{Je sais identifier l'angle dans le milieu incident.}
\checkline{Je sais décider de l'existence d'un rayon réfracté.}
\checkline{Je sais obtenir une image par calcul et par construction.}
\checkline{Je sais démontrer la condition de guidage et le temps de parcours.}
\end{methode}

\Page{Aides progressives}{Des indices pour repartir}{indices}
Ne lire que la première indication utile. Chercher à nouveau avant de passer à la suivante.
\Sub{Exercices}
\begin{description}[style=nextline,leftmargin=0pt,font=\sffamily\bfseries\small]
\item[E04--E06] La normale et la surface font $90^\circ$. Dans la lame, les deux normales sont parallèles ; isoler le triangle qui relie l'épaisseur au trajet interne.
\item[E10--E12] Écrire $1/q=1/f'+1/p$ avec $p=\oa<0$. Ne pas introduire une distance objet positive dans cette écriture.
\item[E14] $\overline{FA}$ est immédiatement connu ; calculer ensuite la distance à partir de $O$.
\item[E15] Poser $x$ et $D-x$ comme deux longueurs positives ; le produit doit valoir $Df'$.
\item[E18] Pour cadrer, imposer $h\,f'/(p-f')\leq h_c$ avec des unités cohérentes.
\item[E20] Les triangles formés après la lentille donnent un rapport de longueurs. Commencer par les plans images.
\item[E23] L'image de $L_1$ devient l'objet de $L_2$. Sa position absolue et sa coordonnée relative à $O_2$ sont différentes.
\item[E24] Dériver par rapport à $p$ en gardant $q$ constant, puis réciproquement. Les incertitudes s'additionnent en variance sous l'hypothèse d'indépendance.
\end{description}
\Sub{Défi D1}
\textbf{Indice 1.} Écrire $1/q=1/f'-1/p$ et utiliser $1/q_0=1/f'-1/p_0$.\par
\textbf{Indice 2.} L'expression dans la valeur absolue devient $q_0(1/p-1/p_0)$.\par
\textbf{Indice 3.} Encadrer $1/p$ puis discuter le signe de sa borne inférieure.
\Sub{Défi D2}
\textbf{Indice 1.} Le trajet le plus lent est la limite de guidage.\par
\textbf{Indice 2.} La durée initiale consomme déjà une partie du budget temporel.\par
\textbf{Indice 3.} Isoler $L$ dans $\tau_0+\kappa L\leq1/(2R)$.
\Sub{Défi D3 et DM}
Pour le microscope, partir de la condition sur l'image finale et remonter vers l'objet. Pour Bessel, relier l'écart des deux racines au discriminant avant de faire les applications numériques.
